%==============================================================================
%  Abstract — keep within ~250–350 words. Mirror the IMRaD logic in miniature:
%  context, problem, theory, methods, findings, contribution.
%==============================================================================
\section*{Abstract}
\addcontentsline{toc}{section}{Abstract}

\noindent
% --- 1. Context (1–2 sentences) -----------------------------------------------
The European Innovation Council (EIC) has positioned itself as the European
Union's flagship instrument for breakthrough innovation, deploying a portfolio
of funding instruments — including the early-stage Pathfinder programme and,
more recently, the Strategic Technologies for Europe Platform (STEP) — that
together signal a shift in how public capital is mobilised for deep-tech and
strategically important technologies.

% --- 2. Problem / research gap (1–2 sentences) --------------------------------
Pathfinder and STEP Scale Up share an institutional home and a single annual
Work Programme, yet operate by visibly different rules: non-dilutive grants
to early-stage research consortia on one side, equity-style catalytic
investment into industrial scale-ups on the other. The thesis asks how
investment logics are constructed, legitimised and contested in the text of
a single, coherent EIC offer that nonetheless presents two very differently
shaped instruments.

% --- 3. Theoretical framing (1–2 sentences) -----------------------------------
Drawing on institutional theory — in particular the literatures on
institutional logics, legitimacy, and field-level isomorphism — this thesis
analyses how investment logics are constructed and contested across the two
instruments.

% --- 4. Method (2–3 sentences) ------------------------------------------------
The study employs a documentary qualitative case design centred on the EIC
Work Programme 2026 (Pathfinder and STEP Scale Up sections). Data are coded
deductively against an institutional-theory codebook (Scott's three pillars,
DiMaggio \& Powell's isomorphism, Thornton et al.'s institutional logics) and
inductively for emergent patterns.

% --- 5. Findings (2–3 sentences) ----------------------------------------------
Three findings stand out. Both instruments move through a
science-to-market-to-state succession along the technology readiness scale,
and the same succession is reproduced within individual challenge chapters.
Where two logics co-occur, the text most often dissolves the tension into a
single instruction rather than holding it open; that dissolution is itself
the cognitive legitimacy work the Work Programme performs. STEP's design
produces the market signal it then claims to respond to, a circular
validation in which the public actor constitutes the field of investible
opportunities, while the Sovereignty Seal recurs across both instruments as
a portable legitimacy artefact that travels between funders.

% --- 6. Contribution (1–2 sentences) ------------------------------------------
The thesis applies institutional theory to a setting the literature has
largely passed over, supranational innovation policy, and develops four
analytical observations from the corpus reading that warrant testing on
broader empirical material. It also offers practical implications for EIC
programme design, applicants, and policy makers concerned with the
alignment of mission-driven funding instruments.

\vspace{0.4cm}
\noindent\textbf{Keywords:} European Innovation Council; EIC Pathfinder;
STEP; institutional theory; institutional logics; investment logics;
innovation policy; deep tech; strategic autonomy.
